
On
définit ensuite un homomorphisme @@M0@@ en associant à @@M1@@ l'élément @@M2@@; cet homomorphisme s'annule sur @@M3@@, et définit par conséquent un homomorphisme @@M4@@. Enfin, on définit un homomorphisme @@M5@@ en associant à @@M6@@ l'élément @@M7@@; cet homomorphisme s'annule sur @@M8@@ et définit donc un homomorphisme @@M9@@.

\textbf{Proposition 4.7.6.}
\begin{enumerate}[label=\roman*)]
\item @@M10@@ \uline{et} @@M11@@ \uline{sont des} @@M12@@-\uline{homomorphismes};
\item @@M13@@ \uline{est} @@M14@@-\uline{linéaire et} @@M15@@ \uline{est} @@M16@@-\uline{linéaire};
\item @@M17@@ \uline{est l'homomorphisme} « \uline{multiplication par} @@M18@@ »;
\item \uline{pour tout} @@M19@@, \uline{et tout} @@M20@@, @@M21@@;
\item @@M22@@.
\end{enumerate}

Il suffit de faire tourner la machine.

Si on en revient à la situation géométrique rappelée plus haut, on voit que @@M23@@ se construit à partir de @@M24@@, @@M25@@, des applications @@M26@@ et @@M27@@ et de l'élément @@M28@@ de @@M29@@ par la méthode précédente, les opérations @@M30@@ sur @@M31@@ ne coïncidant toutefois avec les opérations @@M32@@ définies plus haut sur @@M33@@ qu'à torsion près, tant que la relation (4.3.1) n'est prouvée que modulo torsion.

\textbf{Théorème 4.8.} \uline{Soient} @@M34@@ \uline{une immersion fermée régulière de codimension constante} @@M35@@, @@M36@@ \uline{le schéma obtenu en faisant éclater} @@M37@@, \uline{et} @@M38@@, \uline{de sorte qu'on a le diagramme commutatif} :
@@M39@@
\uline{On suppose que} @@M40@@ \uline{est quasi-compact et a un faisceau ample}. \uline{Alors on a pour tout} @@M41@@ \uline{la suite
exacte} :
